\documentclass{ximera}
% xmPreamble.tex en xmPrintstyle.sty worden automatisch geladen door ximera.cls (v2.7.8);
% NIET nog eens \input{./preamble.tex} of \addPrintStyle{.} doen (geeft 'already defined'-fouten).

\makeatletter % eigen img-map toevoegen aan het zoekpad (werkt standalone én vanuit de xourse)
\g@addto@macro\Ginput@path{{img/}{UW4202/img/}}
\makeatother

\begin{document}
\author{Alexander Holvoet}
\xmtitle{Een hoop zand op een driehoekig grondvlak}{}

\begin{exercise}
    Neem als volgend grondvlak een willekeurige driehoek. Schets eerst wat je verwacht als
    bovenaanzicht, en teken ook de vermoedelijke hoogtelijnen. Giet dan zoveel mogelijk
    zand op de omhoog geplaatste driehoek. Komt het experiment overeen met je schets? Is de
    rusthoek ook hier weer gelijk aan die in de vorige experimenten?
    \begin{oplossing}
        Je krijgt een piramide, zoals je in de foto's hieronder ziet. Die vorm hangt enkel
        af van het grondvlak en niet van waar je begint met zand te gieten. De rusthoek is
        weer dezelfde.
        \begin{image}[12cm]
            
\includegraphics[width=6cm,keepaspectratio]{lzandopdriehoek.jpg}
            \quad
            \includegraphics[width=6cm,keepaspectratio]{lsituatie2.jpg}
        \end{image}
        Een zandhoop op een driehoekig grondvlak: een piramide met de hoogtelijnen erop.
    \end{oplossing}
\end{exercise}

% Opmerking voor de leerkracht: in de bron stond hiernaast ook een lijntekening van de
% hoogtelijnen op de driehoek (Lhoogtelijnendriehoek.png), niet aanwezig in img/.

\begin{exercise}
    Als je ook hier de hoekpunten van de hoogtelijnen verbindt, zie je speciale lijnen
    verschijnen op het bovenaanzicht. Welke lijnen zijn dit?
    \begin{oplossing}
        Het zijn de bissectrices van de hoeken gevormd door twee opeenvolgende zijden van
        de driehoek.
    \end{oplossing}
\end{exercise}

% TODO (afbeelding ontbreekt): \afbeelding[6cm]{Ldriehoekkortsteweg.png}.

\begin{exercise}
    Geef hiervoor een verklaring. Waarmee komen deze lijnen overeen in 3D?
    \begin{oplossing}
        De bissectrices zijn de meetkundige plaats van de punten die even ver liggen van
        twee opeenvolgende zijden van de driehoek. In 3D komen ze overeen met de ribben van
        de piramide. Korrels op de ribben weten niet naar welke kant ze moeten vallen en
        blijven dus liggen.
    \end{oplossing}
\end{exercise}

\begin{exercise}
    We weten dat de drie bissectrices van een driehoek elkaar in \'e\'en punt snijden. Wat
    weet je van dit snijpunt?
    \begin{oplossing}
        Omdat het snijpunt op elk van de bissectrices ligt, ligt het even ver van de drie
        zijden van de driehoek. De loodlijnen op de zijden geven de kortste weg aan voor een
        zandkorrel om naar beneden te rollen.
    \end{oplossing}
\end{exercise}

% TODO (afbeelding ontbreekt): \afbeelding[8cm]{zandopdriehoek.png}.

\begin{exercise}
    Wat betekent het snijpunt van de bissectrices uit het bovenaanzicht voor de zandhoop?
    \begin{oplossing}
        Dit snijpunt is de projectie van de top van de zandhoop. Deze top is het punt dat
        zich even ver van de drie randen van het grondvlak bevindt.
    \end{oplossing}
\end{exercise}

\end{document}
